% TODO(zh): 这节需要你自己的语言组织，下面给出段落级骨架和每段该完成的任务。

% Paragraph 1 -- PINN background and motivation for the specific PDE.
Physics-informed neural networks \citep{raissi2019pinn} solve forward and
inverse partial differential equation (PDE) problems by embedding the
governing equation, boundary conditions, and (for inverse problems)
observational data directly into the training loss of a neural network
surrogate. \todo{One or two sentences motivating the seawater temperature
diffusion equation specifically -- cite the paper you are replicating/
extending (Han et al., 2026, Journal of Oceanology and Limnology) and state
why it is a useful benchmark: closed-form analytical solutions exist for
all three boundary-condition types, which makes ground-truth error
measurable exactly rather than via a numerical reference solution.}

% Paragraph 2 -- QML hype + hybrid quantum-classical PINNs as a response.
\todo{Motivate interest in quantum machine learning for scientific
computing in one or two sentences, then introduce the hybrid
quantum-classical PINN (QCPINN/QPINN) line of work as the object of study.
Keep this short -- the literature survey belongs in Related Work, not here.}

% Paragraph 3 -- the gap / research questions.
Despite growing enthusiasm for hybrid quantum-classical PINNs, we find that
efficiency in this literature is reported almost exclusively via parameter
count or number of training epochs, rarely via wall-clock time -- a
distinction that matters because variational quantum circuits are, in the
near term, executed via classical simulation whose per-step cost scales
with circuit width and depth (see Section~\ref{sec:related}). This raises
two questions that motivate the present study:
\begin{itemize}
  \item[\textbf{RQ1}] On a benchmark PDE with exact analytical solutions,
    does a QCPINN built on a previously proposed circuit ansatz
    \citep{farea2025qcpinn} match classical PINN accuracy once wall-clock
    training time is accounted for, not just parameter count?
  \item[\textbf{RQ2}] If it does not, is the gap attributable to
    training hyperparameters (initialization, learning-rate schedule), or
    to the quantum circuit's structure (expressivity or trainability)?
\end{itemize}

% Paragraph 4 -- contributions, stated as a list.
We make the following contributions:
\begin{itemize}
  \item A controlled classical-vs-quantum comparison across six matched
    scenarios (three boundary-condition types $\times$ forward/inverse),
    with explicit wall-clock accounting alongside accuracy metrics
    (Section~\ref{sec:experiments}).
  \item A two-phase ablation that first rules out training hyperparameters
    as the cause of the quantum branch's underperformance, then isolates
    circuit structure (qubit count, depth) while directly measuring the
    circuit-only gradient norm -- a diagnostic largely absent from prior
    QCPINN evaluations, which infer trainability issues indirectly from
    accuracy curves or rely on purely theoretical gradient-decay bounds
    (Section~\ref{sec:experiments}).
  \item A discussion of why our findings diverge from prior QCPINN results
    obtained with a closely related circuit design, and what that implies
    for how QCPINN efficiency should be reported going forward
    (Section~\ref{sec:discussion}).
\end{itemize}

\todo{Close with a one-sentence roadmap of the paper's structure once the
section numbers/titles are final.}
